\chapter*{Phụ lục}
\addcontentsline{toc}{chapter}{Phụ lục}

\section*{Phụ lục A. Hướng dẫn sử dụng hệ thống}
% TODO

\section*{Phụ lục B. Biểu mẫu khảo sát}
% TODO: nếu đồ án có khảo sát UX tại mục~\ref{sec:danhgiaux} (Ch.5), đính kèm biểu mẫu/câu hỏi khảo sát tại đây.

\section*{Phụ lục C. Mã nguồn chi tiết bổ sung}
% TODO: (nếu cần) các đoạn code không phù hợp trích trong thân báo cáo Ch.4.

% Phụ lục D — Điều khoản dịch vụ. Viết để đọc độc lập: không dùng mã vai trò
% hệ thống, mã yêu cầu chức năng hay tham chiếu chéo vào thân báo cáo.

\section*{Phụ lục D. Điều khoản dịch vụ}
\addcontentsline{toc}{section}{Phụ lục D. Điều khoản dịch vụ}

\noindent\textbf{Điều 1. Phạm vi áp dụng}\\
Điều khoản này áp dụng cho mọi tổ chức, cá nhân sử dụng nền tảng Farmery,
bao gồm nhà cung cấp nông sản, doanh nghiệp mua hàng với số lượng lớn và
người tiêu dùng mua hàng cho nhu cầu cá nhân. Việc đăng ký tài khoản và sử
dụng nền tảng đồng nghĩa với việc chấp nhận toàn bộ các điều khoản dưới đây.

\noindent\textbf{Điều 2. Tài khoản}\\
Mỗi tài khoản chỉ được sử dụng đúng với mục đích đã đăng ký và phạm vi quyền
hạn được Farmery cấp. Người dùng chịu trách nhiệm bảo mật thông tin đăng
nhập và chịu trách nhiệm về toàn bộ hoạt động phát sinh từ tài khoản của
mình; khi nghi ngờ tài khoản bị người khác truy cập trái phép, người dùng
phải thông báo ngay cho Farmery.

\noindent\textbf{Điều 3. Quyền và nghĩa vụ của nhà cung cấp}\\
Nhà cung cấp cam kết cung cấp thông tin sản phẩm, giấy chứng nhận và nhật ký
canh tác trung thực, chính xác; chịu trách nhiệm về chất lượng hàng hóa đúng
như đã mô tả; cập nhật kịp thời số lượng hàng còn lại và giá bán; giao hàng
đúng cam kết về thời gian và điều kiện bảo quản đối với hàng dễ hư hỏng.
Nhà cung cấp có quyền phản hồi công khai đối với đánh giá của người mua.

\noindent\textbf{Điều 4. Quy định về hàng hóa được phép đăng bán}\\
Farmery là nền tảng chuyên biệt cho nông sản. Nhà cung cấp chỉ được đăng bán
hàng hóa thuộc các nhóm sau: nông sản tươi chưa qua chế biến; nông sản đã sơ
chế hoặc bảo quản đơn giản mà không thêm phụ gia; và nông sản chế biến sâu,
với điều kiện bắt buộc phải có bản tự công bố sản phẩm cùng giấy chứng nhận
cơ sở đủ điều kiện an toàn thực phẩm còn hiệu lực.

Không được đăng bán trên nền tảng: thủy sản sống, thịt tươi và động vật
sống; thực phẩm chức năng; thuốc bảo vệ thực vật và vật tư nông nghiệp; mọi
hàng hóa phi thực phẩm; và mọi hàng hóa bị cấm kinh doanh theo quy định của
pháp luật.

Mỗi sản phẩm phải được gắn đúng vào danh mục do Farmery quy định, kèm đầy đủ
thông tin bắt buộc của danh mục đó. Nhà cung cấp đăng sản phẩm sai danh mục
hoặc thuộc nhóm không được phép sẽ bị xử lý theo Điều 8.

\noindent\textbf{Điều 5. Quyền và nghĩa vụ của người mua}\\
Người mua có quyền tra cứu thông tin nguồn gốc hàng hóa trước khi giao dịch,
được bảo vệ theo chính sách đổi trả và hoàn tiền tại Điều 7, và được gửi
đánh giá sau khi đã xác nhận nhận hàng. Người mua có nghĩa vụ thanh toán
đúng hạn, cung cấp thông tin nhận hàng chính xác, và phản hồi trung thực khi
gửi khiếu nại hoặc đánh giá.

\noindent\textbf{Điều 6. Phân định trách nhiệm theo hình thức bán hàng}\\
Nền tảng vận hành hai hình thức bán hàng với trách nhiệm khác nhau:
\begin{itemize}
    \item \textbf{Hàng do nhà cung cấp trực tiếp bán:} nhà cung cấp là bên
    bán trong giao dịch và chịu trách nhiệm về chất lượng, số lượng cũng như
    tiến độ giao hàng. Farmery giữ vai trò trung gian kết nối, bảo đảm minh
    bạch thông tin và hỗ trợ giải quyết tranh chấp.
    \item \textbf{Hàng do Farmery mua lại rồi bán:} Farmery là bên bán trong
    giao dịch và trực tiếp chịu trách nhiệm với người mua về chất lượng hàng
    hóa cũng như việc đổi trả, không phụ thuộc vào trách nhiệm của nhà cung
    cấp ban đầu.
\end{itemize}
Thông tin về bên bán được hiển thị rõ trên trang sản phẩm và trên đơn hàng.

\noindent\textbf{Điều 7. Phí dịch vụ}\\
% TODO: nhóm điền mức phí cụ thể theo phân tích chi phí vận hành thực tế.
Đối với hàng do nhà cung cấp trực tiếp bán, Farmery thu phí hoa hồng trên
mỗi giao dịch thành công. Đối với hàng do Farmery mua lại rồi bán, Farmery
không thu phí của nhà cung cấp mà hưởng chênh lệch giữa giá mua và giá bán.
Ngoài ra, nhà cung cấp có thể tự nguyện sử dụng các dịch vụ trả phí gồm xác
thực chứng nhận nhanh và gói hiển thị nổi bật; vị trí hiển thị mua được luôn
gắn nhãn ``Tài trợ'' để phân biệt với kết quả tìm kiếm thông thường.

Farmery không thu phí đăng ký tài khoản và không thu phí duy trì tài khoản
đối với nhà cung cấp. Việc được ưu tiên hiển thị nhờ điểm uy tín cao là
quyền lợi miễn phí, hoàn toàn độc lập với các dịch vụ trả phí nói trên.

\noindent\textbf{Điều 8. Thanh toán}\\
Giao dịch với người tiêu dùng được thanh toán qua ví điện tử, chuyển khoản
ngân hàng hoặc trả tiền khi nhận hàng.

Đối với giao dịch số lượng lớn có giá trị vượt ngưỡng do Farmery công bố,
tiền của bên mua được một đơn vị trung gian thanh toán được cấp phép giữ
lại. Farmery không trực tiếp giữ tiền của người dùng. Khoản tiền này chỉ
được chuyển cho bên bán sau khi bên mua xác nhận đã nhận đủ hàng đúng cam
kết. Khi phát sinh tranh chấp, khoản tiền được giữ nguyên tại đơn vị trung
gian cho tới khi khiếu nại được giải quyết xong theo Điều 9.

\noindent\textbf{Điều 9. Đổi trả và hoàn tiền}\\
% TODO: nhóm điền khung thời gian gửi khiếu nại và thời gian xử lý tối đa
% sau khi chốt được năng lực vận hành thực tế.
Tùy mức độ và nguyên nhân xác định được, khiếu nại được giải quyết theo một
trong các hướng: hoàn tiền toàn phần, hoàn tiền một phần, đổi sang lô hàng
khác, hoặc từ chối khiếu nại nếu bằng chứng không đầy đủ. Do đặc thù hàng
nông sản dễ hư hỏng, người mua cần gửi khiếu nại kèm hình ảnh minh chứng
trong thời hạn Farmery công bố kể từ khi nhận hàng.

\noindent\textbf{Điều 10. Kiểm duyệt và xử lý vi phạm}\\
% TODO: nhóm điền thời gian xử lý hồ sơ xác minh doanh nghiệp sau khi chốt
% quy trình thẩm định thực tế.
Farmery kiểm duyệt hồ sơ nhà cung cấp, sản phẩm và giấy chứng nhận trước khi
cho hiển thị công khai, đồng thời kiểm tra lại định kỳ trong quá trình hoạt
động. Vi phạm được xử lý theo mức độ tăng dần: nhắc nhở, ẩn sản phẩm, cảnh
cáo kèm trừ điểm uy tín, tạm khóa và cuối cùng là khóa vĩnh viễn tài khoản.

Các hành vi bị xử lý khóa tài khoản ngay lập tức, không qua bước trung gian,
gồm: giả mạo giấy chứng nhận, gian lận trong giao dịch, và thao túng hệ
thống đánh giá.

\noindent\textbf{Điều 11. Sở hữu trí tuệ và dữ liệu}\\
Dữ liệu nhật ký canh tác, hình ảnh sản phẩm và giấy chứng nhận do nhà cung
cấp đưa lên vẫn thuộc quyền sở hữu của nhà cung cấp. Farmery được quyền sử
dụng dữ liệu này để vận hành nền tảng, hiển thị công khai trang thông tin
nguồn gốc sản phẩm, và cải thiện các tính năng hỗ trợ người dùng; với mục
đích cuối, dữ liệu có thể được sử dụng ở dạng ẩn danh hoặc tổng hợp khi
không cần gắn với danh tính cụ thể.

\noindent\textbf{Điều 12. Giới hạn trách nhiệm}\\
Đối với hàng do nhà cung cấp trực tiếp bán, trách nhiệm của Farmery giới hạn
trong phạm vi thông tin đã được xác minh qua quy trình kiểm duyệt; các tranh
chấp về chất lượng nằm ngoài phạm vi xác minh được giải quyết theo Điều 9.
Farmery không tự sản xuất và không tự vận chuyển hàng hóa; việc vận chuyển
do các đơn vị dịch vụ độc lập thực hiện.

\noindent\textbf{Điều 13. Chấm dứt hoặc tạm ngừng dịch vụ}\\
Người dùng có quyền tự nguyện đóng tài khoản bất kỳ lúc nào, với điều kiện
không còn đơn hàng hoặc hợp đồng đang thực hiện dở dang. Farmery có quyền
tạm ngừng hoặc chấm dứt tài khoản theo Điều 10, hoặc khi phát hiện dấu hiệu
gian lận nghiêm trọng. Việc chấm dứt không làm mất đi các nghĩa vụ tài chính
đã phát sinh trước đó, bao gồm các khoản đang được giữ tại đơn vị trung gian
thanh toán theo Điều 8.

\noindent\textbf{Điều 14. Giải quyết tranh chấp}\\
Tranh chấp giữa các bên trước hết được giải quyết qua thương lượng trực tiếp
trên kênh liên lạc của nền tảng. Nếu không đạt được thỏa thuận, các bên có
thể yêu cầu Farmery hòa giải dựa trên dữ liệu giao dịch và thông tin nguồn
gốc đã được hệ thống ghi nhận làm căn cứ khách quan. Nếu vẫn không giải
quyết được, tranh chấp được xử lý theo quy định của pháp luật Việt Nam.

\noindent\textbf{Điều 15. Thay đổi điều khoản}\\
Farmery có quyền cập nhật điều khoản dịch vụ khi cần, có thông báo trước cho
người dùng qua nền tảng hoặc qua địa chỉ liên lạc đã đăng ký. Việc tiếp tục
sử dụng nền tảng sau thời điểm điều khoản mới có hiệu lực được xem là chấp
nhận thay đổi.


\section*{Phụ lục E. Business Model Canvas đầy đủ}
\label{apdx:bmc}
\addcontentsline{toc}{section}{Phụ lục E. Business Model Canvas đầy đủ}

Bảng~\ref{tab:bmc-full} liệt kê toàn bộ chín khối của khung Business Model
Canvas để tiện đối chiếu tổng thể; nội dung theo từng cụm đã được phân tích
tại mục~\ref{sec:businessmodelcanvas}. Ký hiệu \emph{3P} chỉ luồng bán sỉ qua
sàn, \emph{1P} chỉ luồng bán lẻ do Farmery tự thu mua và bán lại; mục không
ghi ký hiệu là dùng chung cho cả hai luồng.

\begingroup
\small
\renewcommand{\arraystretch}{1.2}
\setlength{\tabcolsep}{5pt}
\begin{longtable}{L{3.2cm} L{11.6cm}}
\caption{Business Model Canvas đầy đủ của Farmery}
\label{tab:bmc-full} \\
\toprule
\rowcolor{tblheadbg}
\textbf{Khối} & \textbf{Nội dung} \\
\midrule
\endfirsthead
\toprule
\rowcolor{tblheadbg}
\textbf{Khối} & \textbf{Nội dung} \\
\midrule
\endhead
\bottomrule
\endfoot
Phân khúc khách hàng & Nhà cung cấp là nông dân và hợp tác xã tại vùng trồng thí điểm; khách sỉ là nhà hàng, siêu thị, nhà phân phối (3P); khách lẻ là người tiêu dùng cá nhân trong bán kính giao hàng (1P). \\
\rowcolor{tblaltbg}
Giá trị đề xuất & \emph{Nhà cung cấp:} hai lựa chọn đầu ra thay vì phụ thuộc thương lái --- tự bán sỉ trên sàn để giữ toàn bộ giá bán, hoặc bán đứt cả lô để nhận tiền ngay. \emph{Khách sỉ (3P):} nguồn cung đã xác thực chứng nhận, giá bậc minh bạch theo sản lượng, truy xuất tới từng lô. \emph{Khách lẻ (1P):} hàng đã được phân loại và đóng gói, một đầu mối duy nhất chịu trách nhiệm chất lượng và đổi trả. \\
Quan hệ khách hàng & Tự phục vụ trên web cho khách lẻ; nhân sự phụ trách đồng hành cho khách sỉ; hỗ trợ vùng qua hotline cho nhà cung cấp ít rành công nghệ; đội thu mua làm việc trực tiếp với nhà cung cấp theo mùa vụ (1P). \\
\rowcolor{tblaltbg}
Kênh phân phối & Web và ứng dụng Farmery; mạng xã hội và tối ưu tìm kiếm; hội nông dân, hợp tác xã, chương trình OCOP; đội kinh doanh tiếp cận khách sỉ (3P); đội thu mua đặt tại vùng trồng (1P). \\
Hoạt động chính & Kiểm duyệt nhà cung cấp, sản phẩm và chứng nhận (3P); kết nối cung cầu, xử lý báo giá và hợp đồng bán sỉ (3P); khảo giá và đặt mua theo lô (1P); phân loại, đóng gói, định giá bán lẻ (1P); quản lý tồn kho và xử lý hàng cận hạn (1P); vận hành nền tảng, hệ thống uy tín và các mô-đun AI. \\
\rowcolor{tblaltbg}
Nguồn lực chính & Nền tảng web và hạ tầng cho mô-đun AI; dữ liệu truy xuất nguồn gốc và nhật ký canh tác tích lũy theo thời gian; đội ngũ kiểm duyệt, chăm sóc khách hàng, phụ trách khách sỉ; vốn lưu động để thu mua trước khi bán được hàng (1P); kho sơ chế và đóng gói quy mô nhỏ tại vùng trồng (1P). \\
Đối tác chính & Đơn vị vận chuyển thứ ba (3PL); cổng thanh toán trung gian được cấp phép; cơ quan cấp chứng nhận VietGAP và GlobalGAP; hội nông dân và hợp tác xã địa phương; chương trình OCOP. \\
\rowcolor{tblaltbg}
Dòng doanh thu & Hoa hồng trên mỗi giao dịch bán sỉ thành công (3P); biên lợi nhuận bán lại --- chênh lệch giữa giá thu mua và giá bán lẻ (1P); phí dịch vụ giá trị gia tăng gồm xác thực chứng nhận nhanh và gói hiển thị nổi bật. Không thu phí đăng ký hay duy trì tài khoản nhà cung cấp. \\
Cơ cấu chi phí & Giá vốn hàng mua vào --- khoản chi lớn nhất (1P); hao hụt sau thu hoạch và hàng không bán hết trước hạn (1P); sơ chế, đóng gói, vận chuyển chặng đầu về kho (1P); vận hành nền tảng gồm hạ tầng và đội ngũ kỹ thuật; nhân sự kiểm duyệt, chăm sóc khách hàng, thu mua; marketing; chi phí tích hợp bên thứ ba. \\
\end{longtable}
\endgroup

% Các bảng dài tách khỏi thân báo cáo, được \input từ backmatter/phuluc.tex

\section*{Phụ lục F. Ánh xạ vấn đề, hướng giải quyết và mục tiêu}
\label{apdx:anhxa}
\addcontentsline{toc}{section}{Phụ lục F. Ánh xạ vấn đề, hướng giải quyết và mục tiêu}

Bảng~\ref{tab:vande-muctieu} ánh xạ từng vấn đề nêu tại mục~\ref{sec:vande}
tới hướng giải quyết của Farmery và các mục tiêu cụ thể tương ứng tại
mục~\ref{sec:muctieu}.

\begingroup
\small
\renewcommand{\arraystretch}{1.3}
\setlength{\tabcolsep}{5pt}
\begin{longtable}{L{3.3cm} L{7.2cm} L{2.4cm}}
\caption{Ánh xạ giữa vấn đề, hướng giải quyết và mục tiêu của đồ án}
\label{tab:vande-muctieu} \\
\toprule
\rowcolor{tblheadbg}
\textbf{Vấn đề} & \textbf{Hướng giải quyết của Farmery} & \textbf{Mục tiêu} \\
\midrule
\endfirsthead
\toprule
\rowcolor{tblheadbg}
\textbf{Vấn đề} & \textbf{Hướng giải quyết của Farmery} & \textbf{Mục tiêu} \\
\midrule
\endhead
\bottomrule
\endfoot
\textbf{V1.} Cung dồn theo mùa vụ &
Đặt hàng bán sỉ theo hợp đồng và lịch giao để nhà cung cấp chủ động sản xuất; nền tảng chủ động thu mua theo lô ở luồng bán lẻ, tạo thêm một đầu ra ổn định; mô-đun AI dự báo giá theo mùa vụ và vùng trồng. &
MT3, MT4, MT6 \\
\rowcolor{tblaltbg}
\textbf{V2.} Thiếu kênh tiêu thụ &
Hai luồng tiêu thụ song song trên cùng một tập nguồn cung, mở rộng thị trường ra ngoài phạm vi thương lái địa phương; nhà cung cấp không đủ năng lực bán lẻ vẫn tiêu thụ được hàng qua luồng nền tảng tự thu mua. &
MT3, MT4 \\
\textbf{V3.} Hao hụt, thiếu kho lạnh &
Quy trình vận chuyển riêng cho hàng dễ hư hỏng, ưu tiên giao nhanh qua đối tác vận chuyển thứ ba; quản lý tồn kho theo lô với nguyên tắc xuất trước hàng hết hạn trước; giới hạn bán kính giao hàng lẻ để rút ngắn thời gian hàng nằm trên đường. &
MT4, MT5 \\
\rowcolor{tblaltbg}
\textbf{V4.} Thiếu minh bạch, niềm tin &
Xác thực nhà cung cấp và kiểm duyệt chứng nhận; mã QR truy xuất gắn nhật ký canh tác theo lô; hệ thống đánh giá và điểm uy tín; ở luồng bán lẻ, nền tảng trực tiếp chịu trách nhiệm chất lượng thay vì để người mua tự đòi từng người bán nhỏ lẻ. &
MT1, MT2, MT4, MT7 \\
\end{longtable}
\endgroup

\section*{Phụ lục G. Đối chiếu nông sản với các loại hàng hóa khác}
\label{apdx:dacdiem}
\addcontentsline{toc}{section}{Phụ lục G. Đối chiếu nông sản với các loại hàng hóa khác}

Bảng~\ref{tab:dacdiem-hanghoa} đối chiếu nông sản tươi với hàng công nghiệp
và hàng tiêu dùng đóng gói theo từng tiêu chí, kèm ràng buộc thiết kế mà
mỗi khác biệt đặt ra cho hệ thống. Phần phân tích tóm tắt nằm tại
mục~\ref{subsec:sosanhnongsan}.

\begingroup
\small
\renewcommand{\arraystretch}{1.3}
\setlength{\tabcolsep}{5pt}
\begin{longtable}{L{2.6cm} L{3.4cm} L{3.4cm} L{4.6cm}}
\caption{Đối chiếu đặc điểm nông sản với các loại hàng hóa khác}
\label{tab:dacdiem-hanghoa} \\
\toprule
\rowcolor{tblheadbg}
\textbf{Tiêu chí} & \textbf{Nông sản tươi} & \textbf{Hàng công nghiệp, tiêu dùng đóng gói} & \textbf{Ràng buộc đặt ra cho hệ thống} \\
\midrule
\endfirsthead
\toprule
\rowcolor{tblheadbg}
\textbf{Tiêu chí} & \textbf{Nông sản tươi} & \textbf{Hàng công nghiệp, tiêu dùng đóng gói} & \textbf{Ràng buộc đặt ra cho hệ thống} \\
\midrule
\endhead
\bottomrule
\endfoot
Hạn sử dụng & Vài ngày đến vài tuần & Nhiều tháng đến nhiều năm & Quản lý tồn kho theo lô kèm hạn dùng; xuất trước lô hết hạn trước \\
\rowcolor{tblaltbg}
Nguồn cung & Theo mùa vụ, không điều chỉnh được trong ngắn hạn & Theo kế hoạch sản xuất, co giãn theo cầu & Cần dự báo nhu cầu và giá theo mùa vụ, vùng trồng \\
Biến động giá & Mạnh, theo ngày và theo vụ & Ổn định, thay đổi theo chính sách giá & Giá phải cập nhật thường xuyên; biên lợi nhuận dễ bị bào mòn \\
\rowcolor{tblaltbg}
Tính đồng nhất & Khác nhau giữa các lô cùng giống & Đồng nhất theo mã hàng & Dữ liệu mô tả gắn theo lô, không gắn theo mã sản phẩm \\
Đơn vị bán & Khối lượng, bó, thùng; khác nhau giữa lẻ và sỉ & Theo cái, theo hộp & Hỗ trợ nhiều đơn vị tính và quy đổi giữa chúng \\
\rowcolor{tblaltbg}
Khả năng kiểm chứng chất lượng & Người mua không tự kiểm chứng được, kể cả sau khi dùng & Kiểm chứng được bằng quan sát hoặc dùng thử & Bắt buộc có truy xuất nguồn gốc và chứng nhận bên thứ ba \\
Ràng buộc pháp lý & An toàn thực phẩm, kiểm dịch & Chủ yếu về nhãn mác, tiêu chuẩn kỹ thuật & Kiểm duyệt không thể chỉ dựa vào tự động hóa \\
\rowcolor{tblaltbg}
Hao hụt vận chuyển & 20--40\% với rau quả tại Việt Nam & Gần như không đáng kể & Chi phí hao hụt phải tính vào mô hình kinh doanh \\
Điều kiện bảo quản & Phụ thuộc chuỗi lạnh & Nhiệt độ thường & Phụ thuộc năng lực đối tác vận chuyển, nằm ngoài kiểm soát trực tiếp \\
\end{longtable}
\endgroup

\section*{Phụ lục H. Chứng nhận nông sản phổ biến trên thế giới}
\label{apdx:chungnhan}
\addcontentsline{toc}{section}{Phụ lục H. Chứng nhận nông sản phổ biến trên thế giới}

Bảng~\ref{tab:chungnhan} tổng hợp các chứng nhận và giấy tờ mà một lô nông
sản cần có để ra thị trường, kèm mức độ hỗ trợ trong phạm vi đồ án. Phần
phân tích tóm tắt nằm tại mục~\ref{subsec:chungnhan}.

\begingroup
\small
\renewcommand{\arraystretch}{1.3}
\setlength{\tabcolsep}{5pt}
\begin{longtable}{L{3.1cm} L{4.1cm} L{3.6cm} L{3.2cm}}
\caption{Chứng nhận và giấy tờ nông sản theo thị trường}
\label{tab:chungnhan} \\
\toprule
\rowcolor{tblheadbg}
\textbf{Tên} & \textbf{Phạm vi chứng nhận} & \textbf{Thị trường đích} & \textbf{Farmery hỗ trợ} \\
\midrule
\endfirsthead
\toprule
\rowcolor{tblheadbg}
\textbf{Tên} & \textbf{Phạm vi chứng nhận} & \textbf{Thị trường đích} & \textbf{Farmery hỗ trợ} \\
\midrule
\endhead
\bottomrule
\endfoot
\multicolumn{4}{l}{\textbf{\textit{Thực hành sản xuất}}} \\
VietGAP & Thực hành nông nghiệp tốt, khoảng 70 tiêu chí & Nội địa & Có, kiểm duyệt \\
\rowcolor{tblaltbg}
GlobalG.A.P. & Thực hành nông nghiệp tốt, 252 tiêu chuẩn & Xuất khẩu & Có, kiểm duyệt \\
GRASP & Điều kiện lao động, bổ trợ GlobalG.A.P. & Xuất khẩu & Ghi nhận \\
\rowcolor{tblaltbg}
USDA Organic, EU Organic, JAS Organic & Canh tác hữu cơ & Hoa Kỳ, châu Âu, Nhật Bản & Có, kiểm duyệt \\
\multicolumn{4}{l}{\textbf{\textit{An toàn thực phẩm, chế biến và đóng gói}}} \\
HACCP & Phân tích mối nguy và điểm kiểm soát tới hạn & Nội địa và xuất khẩu & Có, kiểm duyệt \\
\rowcolor{tblaltbg}
ISO 22000, FSSC 22000 & Hệ thống quản lý an toàn thực phẩm & Xuất khẩu & Ghi nhận \\
BRCGS Food Safety, IFS Food & Chuẩn do chuỗi bán lẻ lớn yêu cầu & Anh, châu Âu & Ghi nhận \\
\multicolumn{4}{l}{\textbf{\textit{Thị trường ngách}}} \\
\rowcolor{tblaltbg}
Halal, Kosher & Yêu cầu tôn giáo & Thị trường tương ứng & Ghi nhận \\
Fairtrade, Rainforest Alliance & Thương mại công bằng, bền vững & Châu Âu, Bắc Mỹ & Ghi nhận \\
\multicolumn{4}{l}{\textbf{\textit{Giấy tờ đi theo từng lô hàng}}} \\
\rowcolor{tblaltbg}
Kiểm dịch thực vật & Tình trạng dịch hại của lô hàng & Xuất khẩu & Ngoài phạm vi \\
Giấy chứng nhận xuất xứ & Nguồn gốc xuất xứ, ưu đãi thuế quan & Xuất khẩu & Ngoài phạm vi \\
\rowcolor{tblaltbg}
Phiếu kết quả phân tích & Dư lượng thuốc bảo vệ thực vật, kim loại nặng & Nội địa và xuất khẩu & Ngoài phạm vi \\
Mã số vùng trồng, mã cơ sở đóng gói & Định danh nơi trồng và nơi đóng gói; cấp theo Nghị định 38/2026/NĐ-CP & Xuất khẩu chính ngạch & Ghi nhận \\
\rowcolor{tblaltbg}
Điều kiện an toàn thực phẩm, tự công bố sản phẩm & Điều kiện cơ sở và hồ sơ sản phẩm chế biến & Nội địa & Có, bắt buộc với nhóm chế biến sâu \\
OCOP & Xếp hạng sản phẩm đặc trưng địa phương & Nội địa & Có, kiểm duyệt \\
\end{longtable}
\endgroup


% Phụ lục I — bảng đối sánh giải pháp công nghệ, tách khỏi thân báo cáo.
% Phần kết luận và lý do chính nằm tại mục 5.2.

\section*{Phụ lục I. Đối sánh các giải pháp công nghệ}
\label{apdx:congnghe}
\addcontentsline{toc}{section}{Phụ lục I. Đối sánh các giải pháp công nghệ}

Bảng~\ref{tab:sosanh-congnghe} đối sánh các phương án ứng viên cho từng hạng
mục công nghệ, Bảng~\ref{tab:sosanh-ai} đối sánh các hướng tiếp cận cho từng
mô-đun AI. Bộ tiêu chí đối sánh và kết luận lựa chọn được trình bày tại
mục~\ref{sec:giaiphapcongnghe}.

\begingroup
\small
\renewcommand{\arraystretch}{1.3}
\setlength{\tabcolsep}{5pt}
\begin{longtable}{L{2.5cm} L{3.4cm} L{4.2cm} L{3.9cm}}
\caption{Đối sánh các phương án công nghệ theo hạng mục}
\label{tab:sosanh-congnghe} \\
\toprule
\rowcolor{tblheadbg}
\textbf{Hạng mục} & \textbf{Phương án ứng viên} & \textbf{Nhận xét chính} & \textbf{Kết luận} \\
\midrule
\endfirsthead
\toprule
\rowcolor{tblheadbg}
\textbf{Hạng mục} & \textbf{Phương án ứng viên} & \textbf{Nhận xét chính} & \textbf{Kết luận} \\
\midrule
\endhead
\bottomrule
\endfoot
Giao diện & Next.js (React) & Hỗ trợ sẵn server-side rendering (SSR), cần cho trang truy xuất công khai & \textbf{Chọn} \\
& Nuxt (Vue) & Năng lực tương đương, nhưng hệ sinh thái thư viện thương mại điện tử mỏng hơn & Không chọn \\
\rowcolor{tblaltbg}
& Angular & Cấu trúc chặt chẽ nhưng chi phí học cao, dư thừa so với quy mô đồ án & Không chọn \\
Back-end & NestJS (Node, TypeScript) & Dùng chung ngôn ngữ với giao diện; cấu trúc mô-đun và cơ chế phân quyền có sẵn & \textbf{Chọn} \\
\rowcolor{tblaltbg}
& Spring Boot (Java) & Rất trưởng thành, nhưng thêm một ngôn ngữ thứ hai vào dự án & Không chọn \\
& Django, FastAPI (Python) & Trùng ngôn ngữ với dịch vụ AI, nhưng phần nghiệp vụ lõi cần cấu trúc mô-đun rõ hơn & Không chọn \\
\rowcolor{tblaltbg}
& Laravel (PHP) & Phù hợp thương mại điện tử truyền thống, hệ sinh thái AI hạn chế & Không chọn \\
Kiểu kiến trúc & Modular monolith & Giữ được giao dịch trong một cơ sở dữ liệu, vẫn tách được về sau & \textbf{Chọn} \\
\rowcolor{tblaltbg}
& Monolith phân lớp thuần & Đơn giản nhất nhưng không giữ được đường biên miền nghiệp vụ & Không chọn \\
& Microservices & Chi phí vận hành cao; phá vỡ giao dịch chống bán vượt tồn kho & Không chọn \\
\rowcolor{tblaltbg}
Cơ sở dữ liệu & PostgreSQL & Transaction có row-level locking, dữ liệu không gian, tìm kiếm vector, kiểu JSON có chỉ mục --- gom bốn nhu cầu vào một thành phần & \textbf{Chọn} \\
& MySQL & Đáp ứng giao dịch nhưng thiếu dữ liệu không gian và tìm kiếm vector & Không chọn \\
\rowcolor{tblaltbg}
& MongoDB & Linh hoạt lược đồ nhưng xử lý transaction có row-level locking kém hơn & Không chọn \\
Tìm kiếm & Full-text search của PostgreSQL & Đủ cho quy mô mục tiêu, không phát sinh đồng bộ dữ liệu & \textbf{Chọn} \\
\rowcolor{tblaltbg}
& Elasticsearch, Meilisearch & Mạnh hơn nhiều nhưng thêm một hạ tầng phải đồng bộ và vận hành & Để dành cho giai đoạn sau \\
Lưu trữ tệp & Object storage tương thích S3 & Tách tệp khỏi cơ sở dữ liệu, thuận lợi khi nhân bản máy chủ & \textbf{Chọn} \\
\rowcolor{tblaltbg}
& Lưu trong cơ sở dữ liệu & Làm phình dung lượng, chậm sao lưu & Không chọn \\
& Lưu trên đĩa máy chủ ứng dụng & Cản trở việc nhân bản khi mở rộng & Không chọn \\
\rowcolor{tblaltbg}
Đệm và hàng đợi & Redis & Một thành phần phục vụ cả cache lẫn hàng đợi tác vụ nền & \textbf{Chọn} \\
& Message broker chuyên dụng & Năng lực vượt nhu cầu hiện tại & Để dành cho giai đoạn sau \\
\rowcolor{tblaltbg}
Thanh toán & VNPay (môi trường thử nghiệm) & Tài liệu tiếng Việt đầy đủ, hỗ trợ cơ chế tạm giữ tiền & \textbf{Chọn} \\
Vận chuyển & Giao Hàng Nhanh & Giao diện lập trình công khai, độ phủ tốt tại nông thôn & \textbf{Chọn} \\
\end{longtable}
\endgroup

\begingroup
\small
\renewcommand{\arraystretch}{1.3}
\setlength{\tabcolsep}{5pt}
\begin{longtable}{L{2.5cm} L{3.6cm} L{4.5cm} L{3.4cm}}
\caption{Đối sánh các hướng tiếp cận cho từng mô-đun AI}
\label{tab:sosanh-ai} \\
\toprule
\rowcolor{tblheadbg}
\textbf{Mô-đun} & \textbf{Hướng tiếp cận} & \textbf{Nhận xét chính} & \textbf{Kết luận} \\
\midrule
\endfirsthead
\toprule
\rowcolor{tblheadbg}
\textbf{Mô-đun} & \textbf{Hướng tiếp cận} & \textbf{Nhận xét chính} & \textbf{Kết luận} \\
\midrule
\endhead
\bottomrule
\endfoot
Gợi ý sản phẩm & Collaborative filtering thuần & Không hoạt động khi chưa có lịch sử tương tác & Không chọn \\
\rowcolor{tblaltbg}
& Content-based filtering thuần & Chạy được ngay nhưng gợi ý đơn điệu khi dữ liệu tăng & Không chọn \\
& Mô hình hybrid & Dùng đặc trưng nông sản khi chưa có dữ liệu, chuyển dần sang hành vi & \textbf{Chọn} \\
\rowcolor{tblaltbg}
& Mạng nơ-ron & Đòi hỏi lượng tương tác lớn hơn nhiều so với khả năng thu thập & Không chọn \\
Dự báo giá & Mô hình time series có mùa vụ & Phù hợp chuỗi time series ngắn nhưng có tính mùa vụ rõ; dễ giải thích & \textbf{Chọn} \\
\rowcolor{tblaltbg}
& Gradient boosting trên đặc trưng thời gian & Khai thác được biến ngoại sinh; dùng làm mô hình đối chứng & Chọn làm đối chứng \\
& Mạng nơ-ron hồi quy (RNN, LSTM) & Cần lượng dữ liệu vượt xa khả năng thu thập trong đồ án & Không chọn \\
\rowcolor{tblaltbg}
Trợ lý ảo & Intent classification cổ điển & Cần gán nhãn và bảo trì tập intent thủ công, công sức lớn & Không chọn \\
& LLM dùng trực tiếp & Nguy cơ bịa đặt thông tin về chính sách và sản phẩm & Không chọn \\
\rowcolor{tblaltbg}
& Mô hình ngôn ngữ lớn kèm truy hồi tri thức & Câu trả lời dựa trên tài liệu thật của nền tảng; hiểu tiếng Việt tốt & \textbf{Chọn} \\
Giám sát bất thường & Luật ngưỡng cấu hình được & Kiểm chứng được, minh bạch, không cần dữ liệu huấn luyện & \textbf{Chọn} \\
\rowcolor{tblaltbg}
& Unsupervised learning & Không đủ lưu lượng thật để học phân bố bình thường & Hướng phát triển \\
& Supervised learning & Không có tập nhãn gian lận đã xác nhận để đo chất lượng & Hướng phát triển \\
\end{longtable}
\endgroup


% Phụ lục J — bảng truy vết yêu cầu, tách khỏi thân báo cáo.
% User story và phần dẫn nằm tại mục 4.3.

\section*{Phụ lục J. Truy vết từ nhu cầu người dùng tới yêu cầu}
\label{apdx:yeucau}
\addcontentsline{toc}{section}{Phụ lục J. Truy vết từ nhu cầu người dùng tới yêu cầu}

Bảng~\ref{tab:truyvet-us} truy vết từ từng user story nêu tại
mục~\ref{sec:yeucau} sang các mã yêu cầu tương ứng.
Bảng~\ref{tab:fr-nguoidung} đối chiếu nhóm yêu cầu chức năng với nhóm người
dùng sử dụng trực tiếp.

\begingroup
\small
\renewcommand{\arraystretch}{1.3}
\setlength{\tabcolsep}{5pt}
\begin{longtable}{L{2.6cm} L{8.2cm} L{3.0cm}}
\caption{Truy vết từ user story sang yêu cầu}
\label{tab:truyvet-us} \\
\toprule
\rowcolor{tblheadbg}
\textbf{Nhóm} & \textbf{Nhu cầu phát biểu theo user story} & \textbf{Yêu cầu tương ứng} \\
\midrule
\endfirsthead
\toprule
\rowcolor{tblheadbg}
\textbf{Nhóm} & \textbf{Nhu cầu phát biểu theo user story} & \textbf{Yêu cầu tương ứng} \\
\midrule
\endhead
\bottomrule
\endfoot
Nhà cung cấp & Biết giá thị trường trước khi chốt bán, để không bị ép giá khi vào vụ & FR12.2 \\
\rowcolor{tblaltbg}
& Bán được cả lô ngay khi thu hoạch, không phải lo hàng hỏng trong lúc chờ & FR2.3, quy trình thu mua \\
& Chứng minh sản phẩm đạt chuẩn để bán được giá cao hơn & FR3.1--FR3.3, FR4.2 \\
\rowcolor{tblaltbg}
& Thao tác đơn giản, có người hướng dẫn để không bỏ cuộc & NFR6.1--NFR6.3, FR12.3 \\
& Theo dõi doanh thu và đơn hàng của mình & FR2.2 \\
\rowcolor{tblaltbg}
Khách sỉ & Đặt hàng định kỳ với sản lượng ổn định & FR6.3, FR6.5 \\
& Xác minh được chứng nhận của nhà cung cấp & FR3.2, FR4.2 \\
\rowcolor{tblaltbg}
& Giá giảm theo sản lượng đặt & FR6.1, FR6.4 \\
& Yên tâm rằng tiền chỉ chuyển đi khi đã nhận đủ hàng & FR7.3 \\
\rowcolor{tblaltbg}
Khách lẻ & Biết rau mình mua trồng ở đâu, ai trồng & FR4.2, FR4.5 \\
& Có một nơi chịu trách nhiệm khi hàng không đạt & FR8.3, quy trình bán lẻ \\
\rowcolor{tblaltbg}
& Hàng giao nhanh trong ngày để rau còn tươi & FR8.1, FR8.2, NFR1.1 \\
& Tìm được đúng loại nông sản mình cần & FR4.4, FR4.5, FR12.1 \\
\rowcolor{tblaltbg}
Bộ phận thu mua & So sánh chào giá giữa các nhà cung cấp theo mùa vụ & FR12.2, FR6.1 \\
& Được cảnh báo sớm các lô sắp hết hạn để kịp xả hàng & FR2.3, FR11.2 \\
\rowcolor{tblaltbg}
Quản trị viên & Duyệt nhanh hồ sơ và sản phẩm để không làm nản nhà cung cấp mới & FR10.1, FR3.2, FR4.6 \\
& Thấy bức tranh tổng thể về giao dịch và tồn kho & FR10.4, FR10.5 \\
\rowcolor{tblaltbg}
& Xử lý được sản phẩm đăng sai loại hàng & FR4.7, FR10.2 \\
\end{longtable}
\endgroup

\begingroup
\small
\renewcommand{\arraystretch}{1.25}
\setlength{\tabcolsep}{4pt}
\begin{longtable}{L{4.6cm} >{\centering\arraybackslash}p{2.1cm} >{\centering\arraybackslash}p{1.7cm} >{\centering\arraybackslash}p{1.7cm} >{\centering\arraybackslash}p{1.9cm} >{\centering\arraybackslash}p{1.9cm}}
\caption{Nhóm yêu cầu chức năng theo nhóm người dùng}
\label{tab:fr-nguoidung} \\
\toprule
\rowcolor{tblheadbg}
\textbf{Nhóm yêu cầu} & \textbf{Nhà cung cấp} & \textbf{Khách sỉ} & \textbf{Khách lẻ} & \textbf{Thu mua} & \textbf{Quản trị} \\
\midrule
\endfirsthead
\toprule
\rowcolor{tblheadbg}
\textbf{Nhóm yêu cầu} & \textbf{Nhà cung cấp} & \textbf{Khách sỉ} & \textbf{Khách lẻ} & \textbf{Thu mua} & \textbf{Quản trị} \\
\midrule
\endhead
\bottomrule
\endfoot
FR1: Tài khoản \& phân quyền & \checkmark & \checkmark & \checkmark & \checkmark & \checkmark \\
\rowcolor{tblaltbg} FR2: Hồ sơ nhà cung cấp \& tồn kho & \checkmark & & & \checkmark & \\
FR3: Chứng nhận chất lượng & \checkmark & \checkmark & & & \checkmark \\
\rowcolor{tblaltbg} FR4: Sản phẩm \& truy xuất nguồn gốc & \checkmark & \checkmark & \checkmark & \checkmark & \checkmark \\
FR5: Giao dịch bán lẻ & & & \checkmark & \checkmark & \\
\rowcolor{tblaltbg} FR6: Giao dịch bán sỉ & \checkmark & \checkmark & & & \\
FR7: Thanh toán & \checkmark & \checkmark & \checkmark & \checkmark & \\
\rowcolor{tblaltbg} FR8: Vận chuyển \& khiếu nại & \checkmark & \checkmark & \checkmark & & \checkmark \\
FR9: Tương tác người dùng & \checkmark & \checkmark & \checkmark & & \\
\rowcolor{tblaltbg} FR10: Quản trị hệ thống & & & & & \checkmark \\
FR11: Thông báo & \checkmark & \checkmark & \checkmark & \checkmark & \checkmark \\
\rowcolor{tblaltbg} FR12: Hỗ trợ AI & \checkmark & \checkmark & \checkmark & \checkmark & \\
\end{longtable}
\endgroup


% Phụ lục K — từ điển dữ liệu, tách khỏi thân báo cáo (mục 5.3.5).

\section*{Phụ lục K. Từ điển dữ liệu}
\label{apdx:tudien}
\addcontentsline{toc}{section}{Phụ lục K. Từ điển dữ liệu}


Cơ sở dữ liệu được tổ chức theo tám nhóm thực thể tương ứng với các nhóm chức
năng FR1--FR12 tại mục~\ref{sec:yeucauchucnang}, mỗi nhóm một schema riêng theo
ràng buộc RB3 tại mục~\ref{sec:kientruc}. Quy ước: \textbf{PK} = khóa chính,
\textit{FK} = khóa ngoại.

Phụ lục này đặc tả chi tiết \textbf{bốn nhóm cốt lõi} --- người dùng, sản phẩm
và truy xuất nguồn gốc, giao dịch bán lẻ, giao dịch bán sỉ. Bốn nhóm còn lại
(vận chuyển và khiếu nại, tương tác và thông báo, hỗ trợ AI, quản trị và lưu
vết) có cấu trúc đơn giản hơn và được liệt kê ở mức thực thể tại
mục~\ref{subsec:dacta-csdl}; chi tiết trường của chúng sẽ được hoàn thiện cùng
phần hiện thực tại Chương~\ref{chap:hienthuc}.
 
\paragraph*{Nhóm 1 --- Người dùng \& phân quyền (ứng với FR1)}

\noindent Bảng~\ref{tab:db-nhom1} trình bày đặc tả các thực thể thuộc nhóm này.
 
\begingroup
\small
\renewcommand{\arraystretch}{1.2}
\setlength{\tabcolsep}{5pt}
\begin{longtable}{L{4.9cm} L{3.8cm} L{6.0cm}}
\caption{Đặc tả thực thể nhóm 1: người dùng và phân quyền}
\label{tab:db-nhom1} \\
\toprule
\rowcolor{tblheadbg}
\textbf{Trường} & \textbf{Kiểu / Ràng buộc} & \textbf{Mô tả \& FR liên quan} \\
\midrule
\endfirsthead
\toprule
\rowcolor{tblheadbg}
\textbf{Trường} & \textbf{Kiểu / Ràng buộc} & \textbf{Mô tả \& FR liên quan} \\
\midrule
\endhead
\bottomrule
\endfoot
\multicolumn{3}{@{}l}{\cellcolor{tblaltbg}\textbf{User} (Người dùng)} \\
\texttt{id} & PK & Định danh người dùng. \\
\texttt{full\_name} & & Họ tên hiển thị. \\
\texttt{email} & unique, nullable & Kênh đăng nhập và nhận thông báo. \\
\texttt{phone} & unique, nullable & Kênh đăng nhập và nhận OTP (FR1.2). \\
\texttt{password\_hash} & nullable & Băm bằng Bcrypt/Argon2 (NFR3.3); rỗng khi tài khoản chỉ đăng nhập qua mạng xã hội. \\
\texttt{oauth\_provider} & nullable & Nhà cung cấp OAuth (Google, Facebook...) --- FR1.1. \\
\texttt{oauth\_id} & nullable & Định danh tài khoản mạng xã hội tương ứng --- FR1.1. \\
\texttt{role} & enum & \emph{seller, buyer\_retail, buyer\_wholesale, ops\_sourcing, admin\_moderator, admin\_ops, admin} --- trùng đúng tên vai trò RBAC tại mục~\ref{sec:doituongnguoidung}. \\
\texttt{status} & enum & Hoạt động / khóa. \\
\texttt{created\_at} & datetime & Thời điểm tạo tài khoản. \\
\midrule
\multicolumn{3}{@{}l}{\cellcolor{tblaltbg}\textbf{Address} (Địa chỉ) --- sổ địa chỉ giao/nhận hàng, FR1.5} \\
\texttt{id} & PK & \\
\texttt{user\_id} & FK $\to$ User.id & Chủ sở hữu địa chỉ. \\
\texttt{content} & & Địa chỉ đầy đủ. \\
\texttt{label} & & Nhãn gợi nhớ (nhà riêng, công ty...). \\
\texttt{is\_default} & boolean & Địa chỉ mặc định khi đặt hàng. \\
\midrule
\multicolumn{3}{@{}l}{\cellcolor{tblaltbg}\textbf{Business} (Doanh nghiệp) --- hồ sơ doanh nghiệp thu mua, FR1.3} \\
\texttt{id} & PK & \\
\texttt{owner\_user\_id} & FK $\to$ User.id & Người đứng tên tài khoản doanh nghiệp. \\
\texttt{company\_name} & & \\
\texttt{tax\_code} & unique & Mã số thuế dùng để xác minh. \\
\texttt{business\_license\_file} & & Tệp giấy phép kinh doanh đính kèm. \\
\texttt{verification\_status} & enum & Chờ duyệt / đã duyệt / từ chối. \\
\end{longtable}
\endgroup
Quan hệ: một \emph{User} có nhiều \emph{Address} (1--n); một \emph{Business}
gắn với một \emph{User} đại diện.
 
\paragraph*{Nhóm 2 --- Gian hàng, chứng nhận, sản phẩm \& truy xuất nguồn gốc
(ứng với FR2--FR4)}

\noindent Bảng~\ref{tab:db-nhom2} trình bày đặc tả các thực thể thuộc nhóm này.
 
\begingroup
\small
\renewcommand{\arraystretch}{1.2}
\setlength{\tabcolsep}{5pt}
\begin{longtable}{L{4.9cm} L{3.8cm} L{6.0cm}}
\caption[Đặc tả thực thể nhóm 2: gian hàng, sản phẩm và truy xuất nguồn gốc]{Đặc tả thực thể nhóm 2: gian hàng, chứng nhận, sản phẩm và truy xuất nguồn gốc}
\label{tab:db-nhom2} \\
\toprule
\rowcolor{tblheadbg}
\textbf{Trường} & \textbf{Kiểu / Ràng buộc} & \textbf{Mô tả \& FR liên quan} \\
\midrule
\endfirsthead
\toprule
\rowcolor{tblheadbg}
\textbf{Trường} & \textbf{Kiểu / Ràng buộc} & \textbf{Mô tả \& FR liên quan} \\
\midrule
\endhead
\bottomrule
\endfoot
\multicolumn{3}{@{}l}{\cellcolor{tblaltbg}\textbf{Shop} (Gian hàng) --- FR2.1} \\
\texttt{id} & PK & \\
\texttt{user\_id} & FK $\to$ User.id & Ràng buộc \texttt{role = seller}. \\
\texttt{shop\_name} / \texttt{description} & & Thông tin hiển thị của gian hàng. \\
\texttt{growing\_region} & & Khu vực canh tác công bố. \\
\texttt{ekyc\_status} & enum & Chờ duyệt / đã duyệt / từ chối --- FR1.4. \\
\texttt{reputation\_score} & 0--100 & Giá trị dẫn xuất, tính lại sau mỗi đánh giá hợp lệ theo công thức tại mục~\ref{subsec:danhgianghiepvu}. \\
\texttt{reputation\_updated\_at} & datetime & Lần cập nhật điểm uy tín gần nhất. \\
\midrule
\multicolumn{3}{@{}l}{\cellcolor{tblaltbg}\textbf{Certification} (Chứng nhận) --- FR3.1--FR3.3} \\
\texttt{id} & PK & \\
\texttt{shop\_id} & FK $\to$ Shop.id & \\
\texttt{type} & enum & \emph{VietGAP, GlobalGAP, organic, OCOP, HACCP} --- đúng danh sách tại FR3.1. \\
\texttt{certificate\_number} & & Mã số do cơ quan cấp phát hành. \\
\texttt{issuing\_authority} & & Tổ chức cấp chứng nhận. \\
\texttt{evidence\_file} & & Tệp bản scan chứng nhận. \\
\texttt{issued\_date} / \texttt{expiry\_date} & date & Theo dõi hiệu lực, nhắc gia hạn (FR3.3). \\
\texttt{approval\_status} & enum & Chờ / đã duyệt / từ chối (FR3.2). \\
\texttt{reviewer\_id} & FK $\to$ User.id, nullable & Quản trị viên kiểm duyệt. \\
\midrule
\multicolumn{3}{@{}l}{\cellcolor{tblaltbg}\textbf{Category} (Danh mục) --- phân loại đa cấp, FR4.4} \\
\texttt{id} & PK & \\
\texttt{name} & & Tên danh mục. \\
\texttt{parent\_category\_id} & FK $\to$ Category.id, nullable & Tự tham chiếu để tạo cây danh mục. \\
\midrule
\multicolumn{3}{@{}l}{\cellcolor{tblaltbg}\textbf{Product} (Sản phẩm) --- FR4.1} \\
\texttt{id} & PK & \\
\texttt{shop\_id} & FK $\to$ Shop.id & Gian hàng sở hữu sản phẩm. \\
\texttt{category\_id} & FK $\to$ Category.id & \\
\texttt{name} / \texttt{description} & & Thông tin hiển thị. \\
\texttt{unit} & & kg, tấn, thùng... \\
\texttt{retail\_price} & & Giá bán lẻ B2C. \\
\texttt{status} & enum & Đang bán / ẩn / chờ duyệt. \\
\midrule
\multicolumn{3}{@{}l}{\cellcolor{tblaltbg}\textbf{WholesalePriceTier} (Bảng giá sỉ) --- giá bậc theo MOQ, FR4.1 và FR6.4} \\
\texttt{id} & PK & \\
\texttt{product\_id} & FK $\to$ Product.id & \\
\texttt{min\_quantity} & & Ngưỡng số lượng bắt đầu áp dụng bậc giá. \\
\texttt{unit\_price} & & Đơn giá tương ứng bậc. \\
\midrule
\multicolumn{3}{@{}l}{\cellcolor{tblaltbg}\textbf{Batch} (Lô hàng) --- đơn vị truy xuất nguồn gốc, FR2.3 và FR4.2} \\
\texttt{id} & PK & \\
\texttt{product\_id} & FK $\to$ Product.id & \\
\texttt{batch\_code} & unique & Mã lô hàng. \\
\texttt{harvest\_date} / \texttt{expiry\_date} & date & Cơ sở áp dụng nguyên tắc FEFO (mục~\ref{subsec:khobai}). \\
\texttt{stock\_quantity} & & Tồn kho hiện tại của lô. \\
\texttt{qr\_code} & unique & Mã QR truy xuất gắn với lô. \\
\midrule
\multicolumn{3}{@{}l}{\cellcolor{tblaltbg}\textbf{FarmingLog} (Nhật ký canh tác) --- FR4.2} \\
\texttt{id} & PK & \\
\texttt{batch\_id} & FK $\to$ Batch.id & \\
\texttt{activity\_date} & date & Ngày thực hiện hoạt động canh tác. \\
\texttt{activity} / \texttt{description} & & Gieo trồng, chăm sóc, thu hoạch... \\
\texttt{images} & & Ảnh minh chứng. \\
\texttt{is\_published} & boolean & Đã hiển thị trên trang truy xuất hay chưa. \\
\end{longtable}
\endgroup
Quan hệ: Shop 1--n Certification; Shop 1--n Product; Product 1--n
WholesalePriceTier; Product 1--n Batch; Batch 1--n FarmingLog.
\textbf{Ràng buộc toàn vẹn (NFR9.1):} khi \texttt{Certification.approval\_status
= approved} hoặc \texttt{FarmingLog.is\_published = true}, các trường nội
dung không được phép \texttt{UPDATE} trực tiếp --- mọi chỉnh sửa sau công khai
phải tạo bản ghi mới và giữ liên kết lịch sử (append-only), liên hệ
mục~\ref{subsec:gs1blockchain}.
 
\paragraph*{Nhóm 3 --- Giao dịch bán lẻ B2C (ứng với FR5)}

\noindent Bảng~\ref{tab:db-nhom3} trình bày đặc tả các thực thể thuộc nhóm này.
 
\begingroup
\small
\renewcommand{\arraystretch}{1.2}
\setlength{\tabcolsep}{5pt}
\begin{longtable}{L{4.9cm} L{3.8cm} L{6.0cm}}
\caption{Đặc tả thực thể nhóm 3: giao dịch bán lẻ B2C}
\label{tab:db-nhom3} \\
\toprule
\rowcolor{tblheadbg}
\textbf{Trường} & \textbf{Kiểu / Ràng buộc} & \textbf{Mô tả \& FR liên quan} \\
\midrule
\endfirsthead
\toprule
\rowcolor{tblheadbg}
\textbf{Trường} & \textbf{Kiểu / Ràng buộc} & \textbf{Mô tả \& FR liên quan} \\
\midrule
\endhead
\bottomrule
\endfoot
\multicolumn{3}{@{}l}{\cellcolor{tblaltbg}\textbf{Cart} (Giỏ hàng) \textbf{/ CartItem} (Chi tiết giỏ hàng) --- FR5.1} \\
\texttt{Cart.id} & PK & Mỗi người mua một giỏ hàng. \\
\texttt{Cart.user\_id} & FK $\to$ User.id & \\
\texttt{CartItem.id} & PK & \\
\texttt{cart\_id} & FK $\to$ Cart.id & \\
\texttt{product\_id} & FK $\to$ Product.id & Sản phẩm trong giỏ, có thể thuộc nhiều gian hàng. \\
\texttt{quantity} & & \\
\midrule
\multicolumn{3}{@{}l}{\cellcolor{tblaltbg}\textbf{Voucher} (Mã giảm giá) --- FR5.3} \\
\texttt{id} & PK & \\
\texttt{code} & unique & Mã người mua nhập khi thanh toán. \\
\texttt{scope} & enum & \emph{platform} hoặc \emph{shop}. \\
\texttt{shop\_id} & FK $\to$ Shop.id, nullable & Chỉ dùng khi phạm vi là gian hàng. \\
\texttt{discount\_percent} / \texttt{max\_discount} & & Mức giảm và trần giảm giá. \\
\texttt{remaining\_quantity} / \texttt{expiry\_date} & & Giới hạn phát hành và hiệu lực. \\
\midrule
\multicolumn{3}{@{}l}{\cellcolor{tblaltbg}\textbf{Order} (Đơn hàng chính) --- FR5.2} \\
\texttt{id} & PK & \\
\texttt{user\_id} & FK $\to$ User.id & Người mua. \\
\texttt{total\_amount} / \texttt{created\_at} & & Giá trị và thời điểm đặt hàng. \\
\midrule
\multicolumn{3}{@{}l}{\cellcolor{tblaltbg}\textbf{SubOrder} (Đơn hàng con) --- tách theo gian hàng, FR5.2 và FR5.4} \\
\texttt{id} & PK & \\
\texttt{order\_id} & FK $\to$ Order.id & \\
\texttt{shop\_id} & FK $\to$ Shop.id & Mỗi đơn con thuộc đúng một gian hàng. \\
\texttt{status} & enum & Chờ xác nhận / đang giao / hoàn tất / hủy / hoàn tiền. \\
\texttt{shipping\_fee} / \texttt{total\_amount} & & Tính riêng cho từng đơn con (FR8.2). \\
\midrule
\multicolumn{3}{@{}l}{\cellcolor{tblaltbg}\textbf{OrderItem} (Chi tiết đơn hàng)} \\
\texttt{id} & PK & \\
\texttt{sub\_order\_id} & FK $\to$ SubOrder.id & \\
\texttt{batch\_id} & FK $\to$ Batch.id & Trừ tồn theo đúng lô, giữ liên kết truy xuất nguồn gốc (FR4.3). \\
\texttt{quantity} / \texttt{unit\_price} & & \\
\midrule
\multicolumn{3}{@{}l}{\cellcolor{tblaltbg}\textbf{Payment} (Thanh toán) --- dùng chung B2C/B2B, FR7.1--FR7.3} \\
\texttt{id} & PK & \\
\texttt{order\_id} & FK $\to$ Order.id, nullable & Luồng bán lẻ B2C. \\
\texttt{contract\_id} & FK $\to$ Contract.id, nullable & Luồng bán sỉ B2B. \\
\texttt{installment\_id} & FK $\to$ DeliveryInstallment.id, nullable & Thanh toán theo từng đợt giao. \\
\texttt{method} & enum & COD / chuyển khoản / ví điện tử. \\
\texttt{status} & enum & Trạng thái thanh toán. \\
\texttt{holding\_type} & enum & \emph{standard} hoặc \emph{escrow}. \\
\texttt{escrow\_status} & enum & Chưa ký quỹ / đang giữ / đã giải ngân / hoàn trả --- FR7.3. \\
\texttt{gateway\_transaction\_id} & & Mã đối soát với cổng thanh toán trung gian. \\
\texttt{paid\_at} & datetime & \\
\midrule
\multicolumn{3}{@{}l}{\cellcolor{tblaltbg}\textbf{Review} (Đánh giá) --- FR5.5} \\
\texttt{id} & PK & \\
\texttt{type} & enum & \emph{product} hoặc \emph{shop}. \\
\texttt{sub\_order\_id} & FK $\to$ SubOrder.id & Xác định gian hàng được đánh giá và bảo đảm điều kiện đã mua hàng. \\
\texttt{order\_item\_id} & FK $\to$ OrderItem.id, nullable & Chỉ dùng khi \texttt{type = product}. \\
\texttt{user\_id} & FK $\to$ User.id & Người đánh giá. \\
\texttt{rating} & 1--5 & Điểm đánh giá đơn lẻ. \\
\texttt{content} / \texttt{images} & & Nhận xét và ảnh kèm theo. \\
\texttt{moderation\_status} & enum & Công khai / chờ duyệt / đã ẩn --- đánh giá bị ẩn không tính vào điểm uy tín. \\
\texttt{seller\_reply} & nullable & Phản hồi công khai một lần của nhà cung cấp. \\
\end{longtable}
\endgroup
Quan hệ: một Order phân rã thành nhiều SubOrder (1--n, mỗi
SubOrder gắn đúng một Shop); mỗi SubOrder có nhiều
OrderItem; Payment gắn với Order (1--1 hoặc 1--n nếu cho
thanh toán từng phần); Review luôn tham chiếu một SubOrder đã hoàn tất ---
kèm theo một dòng OrderItem khi đánh giá sản phẩm --- để đảm bảo chỉ
đánh giá sau khi đã mua (ràng buộc nghiệp vụ), mỗi đơn hàng con cho phép tối
đa một đánh giá gian hàng và một đánh giá cho mỗi sản phẩm đã mua.
 
\paragraph*{Nhóm 4 --- Giao dịch bán sỉ B2B (ứng với FR6)}

\noindent Bảng~\ref{tab:db-nhom4} trình bày đặc tả các thực thể thuộc nhóm này.
 
\begingroup
\small
\renewcommand{\arraystretch}{1.2}
\setlength{\tabcolsep}{5pt}
\begin{longtable}{L{4.9cm} L{3.8cm} L{6.0cm}}
\caption{Đặc tả thực thể nhóm 4: giao dịch bán sỉ B2B}
\label{tab:db-nhom4} \\
\toprule
\rowcolor{tblheadbg}
\textbf{Trường} & \textbf{Kiểu / Ràng buộc} & \textbf{Mô tả \& FR liên quan} \\
\midrule
\endfirsthead
\toprule
\rowcolor{tblheadbg}
\textbf{Trường} & \textbf{Kiểu / Ràng buộc} & \textbf{Mô tả \& FR liên quan} \\
\midrule
\endhead
\bottomrule
\endfoot
\multicolumn{3}{@{}l}{\cellcolor{tblaltbg}\textbf{RFQ} (Yêu cầu báo giá) --- FR6.1} \\
\texttt{id} & PK & \\
\texttt{business\_id} & FK $\to$ Business.id & Bên mua. \\
\texttt{shop\_id} & FK $\to$ Shop.id & Bên bán nhận yêu cầu. \\
\texttt{product\_id} & FK $\to$ Product.id & \\
\texttt{requested\_quantity} & & Dùng để đối chiếu MOQ của nhà cung cấp. \\
\texttt{desired\_delivery\_time} & & \\
\texttt{quality\_standard} & & VietGAP/GlobalGAP hoặc yêu cầu riêng. \\
\texttt{status} & enum & Mới / đang đàm phán / chốt / từ chối / hết hạn. \\
\midrule
\multicolumn{3}{@{}l}{\cellcolor{tblaltbg}\textbf{RFQNegotiation} (Lịch sử đàm phán RFQ) --- lịch sử đàm phán, FR6.2} \\
\texttt{id} & PK & \\
\texttt{rfq\_id} & FK $\to$ RFQ.id & \\
\texttt{sender\_id} & FK $\to$ User.id & Bên gửi đề xuất. \\
\texttt{proposed\_price} / \texttt{proposed\_quantity} & & Nội dung thương lượng từng lượt. \\
\texttt{content} / \texttt{sent\_at} & & \\
\midrule
\multicolumn{3}{@{}l}{\cellcolor{tblaltbg}\textbf{Contract} (Hợp đồng) --- FR6.3} \\
\texttt{id} & PK & \\
\texttt{rfq\_id} & FK $\to$ RFQ.id & Một RFQ đã chốt sinh đúng một hợp đồng. \\
\texttt{agreed\_price} / \texttt{total\_quantity} & & Điều khoản thương mại đã thống nhất. \\
\texttt{contract\_file} & & Bản hợp đồng điện tử đã sinh. \\
\texttt{signing\_status} / \texttt{signed\_at} & & Chờ ký / đã ký / hủy. \\
\midrule
\multicolumn{3}{@{}l}{\cellcolor{tblaltbg}\textbf{DeliveryInstallment} (Đợt giao hàng) --- FR6.5} \\
\texttt{id} & PK & \\
\texttt{contract\_id} & FK $\to$ Contract.id & \\
\texttt{quantity} & & Sản lượng của đợt giao. \\
\texttt{planned\_delivery\_date} & date & \\
\texttt{status} & enum & Chờ / đang giao / đã giao. \\
\end{longtable}
\endgroup
Quan hệ: RFQ 1--n RFQNegotiation; một RFQ đã chốt sinh đúng một Contract (1--1);
Contract 1--n DeliveryInstallment; Contract/DeliveryInstallment 1--n Payment
(Nhóm 3) --- bảng Payment dùng chung cho cả hai luồng nhờ bộ khóa ngoại
nullable, đồng thời lưu trạng thái tạm giữ tiền của từng đợt giao theo FR7.3.
 


% Phụ lục L — bảng ngưỡng đo lường, tách khỏi thân báo cáo.

\section*{Phụ lục L. Ngưỡng đo lường và căn cứ}
\label{apdx:nguong}
\addcontentsline{toc}{section}{Phụ lục L. Ngưỡng đo lường và căn cứ}

Bảng~\ref{tab:nguongNFR} trình bày ngưỡng đo lường cho các yêu cầu phi chức
năng tại mục~\ref{sec:yeucauphichucnang}; Bảng~\ref{tab:nguongAI} trình bày
ngưỡng chất lượng mục tiêu cho từng mô-đun AI tại mục~\ref{sec:thietkeAI}.
Cả hai bảng đều được xác lập trước khi hiện thực, dùng làm căn cứ đối chiếu
khi đánh giá kết quả tại Chương~\ref{chap:danhgia}.

\begingroup
\small
\renewcommand{\arraystretch}{1.3}
\setlength{\tabcolsep}{5pt}
\begin{longtable}{L{1.5cm} L{6.0cm} L{5.9cm}}
\caption{Ngưỡng đo lường cụ thể cho các yêu cầu phi chức năng}
\label{tab:nguongNFR} \\
\toprule
\rowcolor{tblheadbg}
\textbf{Mã} & \textbf{Đo lường} & \textbf{Căn cứ} \\
\midrule
\endfirsthead
\toprule
\rowcolor{tblheadbg}
\textbf{Mã} & \textbf{Đo lường} & \textbf{Căn cứ} \\
\midrule
\endhead
\bottomrule
\endfoot
\textbf{NFR1.1} &
Thời gian phản hồi tại percentile 95 (p95) dưới 2 giây với tải truy cập thông thường; số đo thực tế đối chiếu tại mục~\ref{sec:danhgiahieunang}. &
Xuất phát từ nhà cung cấp: họ dùng điện thoại tầm trung với mạng 4G không ổn định ở vùng trồng (persona 1, mục~\ref{sec:doituongnguoidung}), nên mọi độ trễ từ phía máy chủ đều cộng dồn lên độ trễ đường truyền vốn đã cao. Nếu thao tác đăng sản phẩm bị chậm, họ sẽ bỏ dở --- đúng điểm đau ``sợ quy trình rườm rà mất thời gian canh tác''. Chọn 2 giây vì đây là ngưỡng người dùng bắt đầu cảm nhận được là chậm; đo ở p95 thay vì trung bình vì chính nhóm mạng yếu nằm ở phần đuôi phân bố, nếu lấy trung bình thì trải nghiệm của họ bị che mất. \\
\rowcolor{tblaltbg}
\textbf{NFR1.2} &
Kiểm thử với tối thiểu 50 yêu cầu đặt hàng đồng thời trên cùng một lô có tồn kho giới hạn; tỷ lệ đơn hàng gây vượt tồn kho phải bằng 0\%. &
Đây là ràng buộc nghiệp vụ cứng, không có mức ``sai số nhỏ chấp nhận được'' -- oversell lớn hơn 0\% gây thiệt hại tài chính và uy tín trực tiếp cho nhà cung cấp (liên hệ mục~\ref{subsec:danhgianghiepvu}). \\
\textbf{NFR2.1} &
Khi bổ sung một vai trò/mô-đun mới, toàn bộ test case hồi quy (regression) của các module hiện có phải PASS -- 0 test case bị phá vỡ. &
Xuất phát từ nhu cầu vận hành: nền tảng sẽ phải thêm vai trò và phân hệ mới khi mở rộng địa bàn hoặc thêm loại hàng, mà mỗi lần thêm lại làm hỏng chức năng cũ thì người dùng đang quen việc sẽ mất niềm tin. Đồ án không đo trực tiếp được ``chi phí thay đổi kiến trúc'' bằng một con số, nên dùng số test case hồi quy bị phá vỡ làm chỉ số gián tiếp nhưng kiểm chứng được. \\
\rowcolor{tblaltbg}
\textbf{NFR2.2} &
Đáp ứng tối thiểu 100 người dùng thao tác tìm kiếm/đặt hàng đồng thời, thời gian phản hồi trung vị dưới 3 giây, tỷ lệ lỗi dưới 1\%, đo bằng công cụ kiểm thử tải (mục~\ref{sec:danhgiahieunang}). &
Xuất phát từ đặc thù hành vi mua nông sản: nhu cầu không trải đều trong ngày mà dồn vào khung sáng sớm và chiều tối, và dồn mạnh hơn nữa vào lúc mở bán một lô hàng mới --- vì hàng tươi có số lượng giới hạn nên khách lẻ có xu hướng vào cùng lúc để tranh mua. Mốc 100 người đồng thời ước lượng cho tình huống dồn tải đó ở địa bàn thí điểm, không phải quy mô của Shopee hay Lazada đã dẫn ở Chương~\ref{chap:lienquan}. \\
\textbf{NFR3.1} &
100\% endpoint API thao tác trên dữ liệu/tài nguyên nhạy cảm được kiểm soát quyền; gọi trực tiếp các endpoint đó với tài khoản không đủ quyền phải trả về lỗi 403 trong mọi trường hợp. &
Xuất phát từ nhu cầu của khách sỉ và nhà cung cấp: khách sỉ cần chắc chắn bảng giá riêng và nội dung đàm phán của mình không lọt sang đối thủ, còn nhà cung cấp cần chắc chắn không ai sửa được nhật ký canh tác hay giá bán của họ. Một endpoint bị bỏ sót đã đủ phá vỡ cam kết đó, nên ngưỡng chấp nhận là 100\%, không có mức trung gian. \\
\rowcolor{tblaltbg}
\textbf{NFR3.2} &
Rà soát cơ sở dữ liệu xác nhận không có trường dữ liệu định danh nào lưu dạng plaintext. &
Đây là dữ liệu cá nhân thuộc phạm vi Nghị định 13/2023/NĐ-CP (NFR8.2); rò rỉ dù một bản ghi cũng là vi phạm nên yêu cầu là tuyệt đối, không đặt ngưỡng tỷ lệ. \\
\textbf{NFR3.3} &
Bcrypt với hệ số chi phí (cost factor) tối thiểu 10, hoặc Argon2id với tham số khuyến nghị tương đương. &
Cân bằng giữa hai nhu cầu đối lập của cùng một nhóm người dùng: nhà cung cấp đăng nhập trên thiết bị yếu nên không chịu được thời gian băm quá lâu, nhưng cũng chính họ là nhóm dễ đặt mật khẩu đơn giản nhất nên cần chi phí băm đủ lớn để cản trở tấn công dò mật khẩu ngoại tuyến. Cost factor 10 là mức khuyến nghị phổ biến đáp ứng được cả hai. \\
\rowcolor{tblaltbg}
\textbf{NFR3.5} &
Khóa tạm thời tài khoản trong 30 phút sau tối đa 5 lần đăng nhập/xác thực sai liên tiếp. &
Thống nhất với ngưỡng đã áp dụng cho luồng xác thực OTP tại mục~\ref{sec:quytrinhnghiepvu} (Đăng ký và xác thực nhà cung cấp), vừa hạn chế dò mật khẩu tự động vừa không gây phiền hà cho người dùng nhập nhầm. \\
\textbf{NFR4.1} &
0 trường hợp sai lệch số liệu tồn kho/số dư phát hiện được qua kiểm thử giao dịch đồng thời (liên hệ NFR1.2). &
Đây là yêu cầu toàn vẹn dữ liệu tài chính/tồn kho, không tồn tại mức ``sai số nhỏ chấp nhận được''. \\
\rowcolor{tblaltbg}
\textbf{NFR4.2} &
Sao lưu tự động tối thiểu 1 lần/ngày, mục tiêu điểm khôi phục (RPO) không quá 24 giờ; thử khôi phục từ bản sao lưu ít nhất một lần trong giai đoạn đánh giá (mục~\ref{sec:danhgiahieunang}). &
Xuất phát từ thiệt hại thực tế nếu mất dữ liệu: nhật ký canh tác là thứ nhà cung cấp phải ghi ngay tại ruộng và gần như không thể khôi phục lại từ trí nhớ, còn dữ liệu tồn kho của luồng 1P mà sai thì kéo theo bán vượt và lỗ trực tiếp. Mất tối đa một ngày dữ liệu là mức thiệt hại còn khắc phục được bằng đối chiếu thủ công; backup liên tục (RPO gần 0) đòi hỏi hạ tầng vượt phạm vi nguồn lực đồ án. \\
\textbf{NFR5.2} &
100\% endpoint thuộc các nhóm FR1--FR12 (mục~\ref{sec:yeucauchucnang}) được document hóa theo OpenAPI/Swagger. &
Xuất phát từ nhu cầu của đội vận hành nội bộ: bộ phận thu mua và quản trị viên phải tra cứu và đối soát dữ liệu hằng ngày, endpoint không có tài liệu là chỗ họ phải hỏi lại đội kỹ thuật. Dùng chính danh sách FRx.x đã đặc tả làm checklist đối chiếu, tránh bỏ sót các endpoint ít được chú ý nhưng vẫn được sử dụng thường xuyên. \\
\rowcolor{tblaltbg}
\textbf{NFR6.1--6.3} &
Trong usability test với đại diện nhóm nông dân/HTX (mục~\ref{sec:danhgiaux}), task completion rate cho luồng đăng bán sản phẩm và ghi nhật ký canh tác đạt tối thiểu 80\%, không cần người quan sát hỗ trợ trực tiếp. &
Xuất phát từ điểm đau đã xác định của nhà cung cấp: họ ngại biểu mẫu phần mềm phức tạp và lo tốn thời gian canh tác (persona 1, mục~\ref{sec:doituongnguoidung}), nên chỉ cần một lần thất bại là họ bỏ hẳn, quay về bán cho thương lái. Chọn hai tác vụ đăng sản phẩm và ghi nhật ký canh tác làm phép đo vì đó là việc họ phải làm lặp đi lặp lại; ngưỡng 80\% nghĩa là cứ năm người thì nhiều nhất một người cần trợ giúp --- mức mà đội hỗ trợ vùng còn kham được. \\
\textbf{NFR7.1} &
100\% tích hợp với bên thứ ba (đơn vị vận chuyển 3PL, cổng thanh toán) thực hiện qua REST API có tài liệu; mỗi lần gọi có thời gian chờ tối đa 5 giây và có cơ chế thử lại khi gặp lỗi tạm thời. &
Hệ thống phụ thuộc vào đối tác bên ngoài mà nhóm không kiểm soát được chất lượng dịch vụ, nên phải chặn thời gian chờ để một API chậm không kéo theo cả luồng đặt hàng vượt ngưỡng NFR1.1. \\
\rowcolor{tblaltbg}
\textbf{NFR7.2} &
Giao diện hiển thị đúng bố cục (không vỡ layout, không cuộn ngang) trên tối thiểu 3 kích thước màn hình đại diện (điện thoại, máy tính bảng, máy tính để bàn) và 2 trình duyệt phổ biến nhất tại Việt Nam. &
Nhóm nông dân/HTX và phần lớn người tiêu dùng cá nhân được xác định truy cập chủ yếu qua điện thoại (mục~\ref{sec:doituongnguoidung}), nên hiển thị tốt trên di động là điều kiện bắt buộc phải kiểm chứng, không chỉ là tính năng ``hỗ trợ thêm''. \\
\textbf{NFR9.1} &
0 bản ghi đã công khai bị chỉnh sửa/xóa trực tiếp -- mọi thay đổi sau khi công khai chỉ qua cơ chế correction có ghi log (bản ghi mới kèm lý do, không ghi đè); kiểm tra qua rà soát log thao tác. &
Mục tiêu là tính bất biến sau khi công khai nên chỉ tiêu đo phải tuyệt đối, nhất quán với lý do chọn giải pháp cơ sở dữ liệu tập trung có kiểm soát thay vì blockchain tại mục~\ref{subsec:gs1blockchain}. \\
\end{longtable}
\endgroup

\begingroup
\small
\renewcommand{\arraystretch}{1.3}
\setlength{\tabcolsep}{5pt}
\begin{longtable}{L{2.2cm} L{5.3cm} L{5.7cm}}
\caption{Ngưỡng chất lượng mục tiêu cho từng mô-đun AI}
\label{tab:nguongAI} \\
\toprule
\rowcolor{tblheadbg}
\textbf{Mô-đun} & \textbf{Chỉ số / ngưỡng mục tiêu} & \textbf{Căn cứ} \\
\midrule
\endfirsthead
\toprule
\rowcolor{tblheadbg}
\textbf{Mô-đun} & \textbf{Chỉ số / ngưỡng mục tiêu} & \textbf{Căn cứ} \\
\midrule
\endhead
\bottomrule
\endfoot
\textbf{Gợi ý sản phẩm} &
Click-through rate (CTR) vào danh sách gợi ý cá nhân hóa cao hơn tối thiểu 20\% so với danh sách baseline không cá nhân hóa (sản phẩm mới nhất/bán chạy chung), đo qua thử nghiệm A/B hoặc so sánh trước--sau khi bật tính năng. &
Đồ án không có đủ dữ liệu lịch sử lớn để đặt mục tiêu precision/recall tuyệt đối như các hệ thống gợi ý thương mại đã vận hành lâu năm; đo uplift CTR so với một baseline không cá nhân hóa khả thi hơn và vẫn phản ánh đúng giá trị thực tế mà tính năng cần tạo ra. \\
\rowcolor{tblaltbg}
\textbf{Dự báo giá} &
MAPE (Mean Absolute Percentage Error) dưới 20\% trên tập kiểm tra tách theo thời gian (time-based split), tương đương độ chính xác từ khoảng 80\% trở lên. Đo trên dữ liệu giá nông sản công khai, không phải dữ liệu giao dịch của nền tảng. &
Đặt thấp hơn một chút so với mức khoảng 82\% Farm2Market (2025) công bố cho bài toán dự báo giá nông sản tương tự \parencite{farm2market_2025}, do đây là công trình nghiên cứu học thuật có thể đã tối ưu trên bộ dữ liệu chọn lọc; 80\% là mục tiêu khả thi, kiểm chứng được trong phạm vi đồ án và vẫn đủ hữu ích để hỗ trợ ra quyết định. \\
\end{longtable}
\endgroup

